\section{Conclusiones}
\label{sec:conclusiones}

\begin{itemize}
  \item El sistema de adquisición guarda el ECG como tensión en voltios, en un CSV de una sola columna, sin base de tiempo ni metadatos. La frecuencia de muestreo, \SI{500}{\hertz}, tuvo que deducirse del firmware y verificarse con la duración de los registros y con la estructura de paquetes de 100 muestras. Comparado con los formatos clínicos y de investigación, es un formato mínimo que obliga a documentar por fuera toda la información necesaria para interpretar la señal.
  \item La señal necesita acondicionamiento: está montada sobre \SI{1.65}{\volt}, tiene deriva de la línea de base, el registro post ejercicio satura y el comienzo de cada archivo es discontinuo. Una resta de la mediana, un pasaaltos de \SI{0.5}{\hertz} de fase cero y una máscara que excluye el inicio y las saturaciones resuelven estos problemas sin desplazar los picos R. No hizo falta filtrar la red eléctrica.
  \item La versión del algoritmo de Pan-Tompkins con umbral fijo detectó todos los latidos de las zonas válidas, sin intervalos atípicos. Referir el umbral al percentil 99 de la señal integrada en lugar de a su máximo fue clave para que un único artefacto no hiciera perder un tercio de los latidos del registro post ejercicio. Al armar el tacograma también fue necesario descartar los intervalos que contienen una zona inválida, y no solo los que tienen un extremo en ella.
  \item Después del ejercicio la frecuencia cardíaca media subió de 84 a \SI{131}{lpm} y bajó \SI{26}{lpm} en tres minutos de recuperación, más rápido en el primero. La variabilidad latido a latido cayó más de cinco veces (RMSSD de 48 a \SI{8.6}{\milli\second}), y no solo por el acortamiento de los intervalos.
  \item El periodograma de Lomb, que trabaja directamente sobre los instantes de los latidos y tolera los huecos del tacograma, mostró una reducción de entre 50 y 100 veces en las potencias LF y HF. El análisis por minuto mostró que la banda HF casi desaparece en el primer minuto de recuperación, con el pico respiratorio corrido a \SI{0.36}{\hertz}, y reaparece en los minutos siguientes.
  \item Los resultados se corresponden con la retirada vagal durante el esfuerzo y su reactivación gradual en la recuperación. El cociente LF/HF del reposo (4,7) es poco confiable, porque el registro es más corto que lo recomendado para estimar la banda LF. Con un único voluntario no fue posible comparar entre sujetos; registros más largos, un segundo voluntario y el registro simultáneo de la respiración permitirían confirmar estas conclusiones.
\end{itemize}
